\textbf{Objective:} To predict 30-day readmission and escalation of
glucose-lowering treatment at the next hospital stay from real-world
longitudinal inpatient records of people with diabetes, and to test whether
knowledge-aware sequence models add value over strong gradient boosting under
rigorous validation.

\textbf{Methods:} We used the public Diabetes 130-US Hospitals database
(\NEnc{} encounters of \NPat{} patients after exclusions; \NRepeat{} patients
with repeat stays). We propose the Temporal Knowledge-Gated Network (TKGN),
which encodes each stay with a hybrid knowledge graph (ICD-9-CM hierarchy,
drug classes and a training-only co-morbidity graph), summarises earlier
stays and fuses them with the current stay through history-reliability
gates, and TKGN-B, which trains TKGN as a residual on an out-of-fold
LightGBM logit. Ten models, including tuned logistic regression, random
forest, XGBoost, LightGBM, GRU, RETAIN and Transformer baselines, were
compared on five patient-disjoint splits and a prospective temporal split,
with patient-clustered bootstrap confidence intervals, DeLong tests,
calibration, ablation, learning-curve, explanation and subgroup analyses.

\textbf{Results:} Readmission and escalation prevalence were
\RateReadmit\% and \RateEsc\%. TKGN-B achieved the highest mean AUROC for
readmission (\AucReadmitGrpTKGNB{} versus \AucReadmitGrpLGBM{} for LightGBM)
and escalation (\AucEscGrpTKGNB{} versus \AucEscGrpLGBM{}), outperforming
LightGBM in \WinsReadmitLGBM{} and \WinsEscLGBM{} repeats, with good
calibration (ECE \EceReadmitGrpTKGNB{} and \EceEscGrpTKGNB{}). It was
significantly better than all deep sequence baselines for readmission and
the best model on the temporal escalation split (\AucEscTmpTKGNB{} versus
\AucEscTmpLGBM{}, $p$ \PDelongTmpEscLGBM{}); differences to gradient boosting on
patient-disjoint splits were small and not statistically significant, and a
random forest was most robust on the temporal readmission split.
Earlier stays were the most valuable model input, whereas knowledge-graph
code sharing did not improve discrimination, even with 5--25\% of the
training data.
Test AUROC differed by at most \GapReadmitRace{} across adequately sized
racial groups and \GapReadmitGender{} between sexes, and most across age
bands (\GapReadmitAge{}).

\textbf{Conclusion:} Anchoring a knowledge-gated sequence model on gradient
boosting yields calibrated risk estimates that match or slightly exceed tuned
gradient boosting and are clearly better than standard deep sequence models
on real diabetes inpatient data. Gains over boosting are modest,
underscoring the need for patient-disjoint and temporal validation when new
architectures are proposed.
